\section*{4 Low-Depth Sorting of Logarithmic-Length Keys (author source pp.11:5--11:9)}
We first recall what it means to sort with respect to a partial ordering.
\begin{originalresult}{Definition 4}
Let $\Sigma$ be a finite set. We say that $x,y\in\Sigma^n$ are reorderings of each other if for some permutation $\pi:[n]\to[n]$, it holds that $x_i=y_{\pi(i)}$, for all $i\in[n]$.
\end{originalresult}
\begin{originalresult}{Definition 5}
Let $\Sigma$ be a finite set with a strict partial ordering $\prec$. A Boolean circuit is said to sort $\Sigma^n$ with respect to $\prec$ if on any input $x\in\Sigma^n$, it outputs a reordering $y$ of $x$ such that for any $i,j\in[n]$, $y_i\prec y_j\implies i<j$.
\end{originalresult}
In particular, we will focus on a partial ordering that compares integers by their $k$ most significant bits. That is, for any $m\in\ZZ^+$ and any $x,y\in\bits^m$, we write $x\prec_k y$ to denote that $(x_1,\ldots,x_k)$ lexicographically precedes $(y_1,\ldots,y_k)$.
\begin{originalresult}{Lemma 6}
For every $m=m(n)$ and $k\in[m]$, there is a log-space uniform Boolean circuit of size $O(n\cdot m\cdot\log n)$ and depth $O(k+\log n)$ for sorting $(\bits^m)^n$ with respect to $\prec_k$.
\end{originalresult}
We are mainly interested in the setting $k=\Theta(\log n)$, in which case the constructed Boolean circuit has size $O(n\cdot m\cdot\log n)$ and depth $O(\log n)$.
To the best of our knowledge Lemma 6 was not previously known. The straight-forward circuit based on AKS networks yields circuits with matching size $O(n\cdot m\cdot\log n)$ but depth $O(\log n\cdot\log k)$. Kospanov [17] obtained circuits with better depth $O(\log n+\log k)$ but much larger size $O(n^2\cdot k)$. A recent work of [18] obtains size $O(n\cdot k^2)$ and depth $O(\log n+k\log k)$, which is better for some values of $k=o(\log n)$, but worse in our regime of $k=\Theta(\log n)$. The work of Lin and Shi [22, Theorem 1.1] similarly constructs circuits with size $O(n\cdot m\cdot\log n)$ and depth $O(\log n)$ when $n>2^{4k+7}$, but not for all values of $k=\Theta(\log n)$.
\subsection*{4.1 Sorting Networks}
Our sorting circuits rely on sorting networks, which are generic constructions of $n$-input sorting circuits from 2-input sorting circuits.
\begin{originalresult}{Definition 7 (Sorting Networks)}
An $n$-input sorting network is a circuit $C$ with $n$ inputs and $n$ outputs, and with all gates having fan-in and fan-out two, such that for any set $\Sigma$ with strict partial ordering $\prec$, if each gate of $C$ is replaced by a circuit that sorts $\Sigma^2$ with respect to $\prec$, then the resulting circuit sorts $\Sigma^n$ with respect to $\prec$.
\end{originalresult}
For the best asymptotic size and depth, we use the celebrated sorting network of Ajtai, Koml\'os, and Szemer\'edi [1].
\begin{originalresult}{Lemma 8 ([1])}
There is a (log-space uniform) $n$-input sorting network with size $O(n\log n)$ and depth $O(\log n)$.
\end{originalresult}
\subsection*{4.2 From Sorting Networks to Boolean Circuits}
We prove Lemma 6 by combining the following lemma with Lemma 8.
\begin{originalresult}{Lemma 9}
Suppose that there is a log-space uniform sorting network on $n$ elements with size $S=S(n)$ and depth $D=D(n)$. Then, for all $m=m(n)$ and $k\in[m]$, there is a log-space uniform Boolean circuit with size $O(S\cdot m)$ and depth $O(D+k)$ for sorting $(\bits^m)^n$ with respect to $\prec_k$.
\end{originalresult}
\begin{proof}
Starting with an $n$-input sorting network $A$ with size $S$ and depth $D$, we obtain a Boolean circuit $C$ from $A$ by replacing each gate with a Boolean circuit $\Sort^{(2)}_{k,m}$ that sorts $(\bits^m)^2$ with respect to $\prec_k$. With this approach it is easy to make $C$ have size $O(S\cdot m)$---since $\Sort^{(2)}_{k,m}$ merely needs to have size $O(m)$.
Depth seems to pose more of a difficulty. Locality considerations imply that $\Sort^{(2)}_{k,m}$ must have depth at least $\Omega(\log k)$ (e.g. the least significant bit of either output of $\Sort^{(2)}_{k,m}(x,y)$ depends on all of the $k$ most significant bits of $x$ and $y$). Thus, it might seem that this approach dooms $C$ to have depth $\Omega(D\cdot\log k)$.
We circumvent this difficulty by noting that although input-to-output paths in $C$ are concatenations of $D$ input-to-output paths in $\Sort^{(2)}_{k,m}$, these $D$ paths cannot all be worst-case. In particular, the $i$th output bit of one copy of $\Sort^{(2)}_{k,m}$ is only connected to the $j$th input bit of another copy if $i\equiv j\pmod m$ (see Figure 1).
\begin{figure}[ht]\centering\includegraphics[page=7,width=.8\textwidth,trim=90 550 105 160,clip]{multiselection-original.pdf}\caption{Original Figure 1: The four different ways that we can connect one copy of $\Sort^{(2)}_{k,n}$ to another. In this depiction, a copy of $\Sort^{(2)}_{k,n}$ is depicted by a rectangle, with inputs incident to the bottom edge and output incident to the top. The left half corresponds to the first input/output, and the right half to the second.}\end{figure}
We leverage this with the following lemma, whose proof we defer momentarily.
\begin{originalresult}{Lemma 10}
For any $m=m(n)$ and $k\in[m]$, there is a log-space uniform Boolean circuit $\Sort^{(2)}_{k,m}$ with size $O(m)$ that sorts $(\bits^m)^2$ with respect to $\prec_k$ and has the following additional property:
For all $\hat i,\hat j\in[m]$ and $i,j\in[2m]$ with $i\equiv\hat i\pmod m$ and $j\equiv\hat j\pmod m$, every path in $\Sort^{(2)}_{k,m}$ from the $i$th input vertex to the $j$th output vertex has length $O(\min(j,k)-\min(i,k)+1)$. In particular there is no such path if $k>i>j$.
\end{originalresult}
Assuming Lemma 10, consider any input-to-output path $p$ in $C$. We write $p$ as a composition of paths $p_1\circ\cdots\circ p_D$, where each $p_\ell$ is an input-to-output path of some copy of $\Sort^{(2)}_{k,m}$. Define $i_\ell,j_\ell\in[2m]$ so that $p_\ell$ starts at the $i_\ell$th input and ends at the $j_\ell$th output of that copy, and define $\hat i_\ell,\hat j_\ell\in[m]$ such that $i_\ell\equiv\hat i_\ell\pmod m$ and $j_\ell\equiv\hat j_\ell\pmod m$. By the observation above, we have $\hat j_\ell=\hat i_{\ell+1}$ for all $\ell\in[D-1]$, and by Lemma 10, we have for some constant $L$ that $|p_\ell|\le L\cdot(\min(\hat j_\ell,k)-\min(\hat i_\ell,k)+1)$.
Putting this together, we bound
\begin{align*}
|p|&=\sum_{\ell=1}^D|p_\ell|\\
&\le\sum_{\ell=1}^D L\cdot(\min(\hat j_\ell,k)-\min(\hat i_\ell,k)+1)\\
&=L\cdot(\min(\hat j_D,k)-\min(\hat i_1,k)+D)\quad\text{(telescoping sum because $\hat j_\ell=\hat i_{\ell+1}$)}\\
&\le L\cdot(k+D)=O(k+D).
\end{align*}
\end{proof}
It remains to prove Lemma 10.
\begin{proof}[Proof of Lemma 10]
We use the straightforward construction that compares the elements bit-by-bit starting with their most significant bits. In more detail, we will construct a circuit $\widetilde{\Sort}^{(2)}_{k,m}:\{\prec,?,\succ\}\times\bits^m\times\bits^m$ that has the functionality
\begin{align*}
\widetilde{\Sort}^{(2)}_{k,m}(?,x,y)&=\begin{cases}(x,y)&\text{if }x\prec y,\\(y,x)&\text{otherwise},\end{cases}\\
\widetilde{\Sort}^{(2)}_{k,m}(\prec,x,y)&=(x,y),\\
\widetilde{\Sort}^{(2)}_{k,m}(\succ,x,y)&=(y,x)
\end{align*}
and we define $\Sort^{(2)}_{k,m}=\widetilde{\Sort}^{(2)}_{k,m}(?,\cdot,\cdot)$. This clearly gives the desired functionality. The first input symbol aids in recursively constructing the circuit for $\widetilde{\Sort}^{(2)}_{k,m}$, and is meant to indicate whether the global decision on the comparison between $x$ and $y$ is (1) still undecided (?), (2) decided toward $x$ or (3) decided toward $y$.
Our construction of $\widetilde{\Sort}^{(2)}_{k,m}$ is recursive. If $k=0$, then we define $\widetilde{\Sort}^{(2)}_{k,m}$ to be the identity function. For $k>0$, $\widetilde{\Sort}^{(2)}_{k,m}$ takes input $(c,x,y)$, it first computes
\[(c',a_1,b_1)=\begin{cases}(\prec,x_1,y_1)&\text{if }c=\prec\text{ or }(c=?\text{ and }x_1<y_1),\\(?,x_1,y_1)&\text{if }c=?\text{ and }x_1=y_1,\\(\succ,y_1,x_1)&\text{if }c=\succ\text{ or }(c=?\text{ and }x_1>y_1).\end{cases}\]
Then, with $x'$ and $y'$ denoting $(x_2,\ldots,x_m)$ and $(y_2,\ldots,y_m)$, it computes
\[(a',b')=\widetilde{\Sort}^{(2)}_{k-1,m-1}(c',x',y').\]
Finally it outputs $a_1\circ a',b_1\circ b')$, where $\circ$ denotes string concatenation. It is easy to check the correctness of this construction.
\paragraph{Circuit Size.} Let $S(k,m)$ denote the size of $\widetilde{\Sort}^{(2)}_{k,m}$. For $k=0$ we clearly have $S(k,m)=O(m)$. For larger $k$, the recursive construction gives $S(k,m)=S(k-1,m-1)+O(1)$, so by induction $S(k,m)=O(m)$ for all $k\le m$.
\paragraph{Input-to-Output Path Lengths.} We now bound the length of different input-to-output paths in $\widetilde{\Sort}^{(2)}_{k,m}$. As in the construction above, denote the input to $\widetilde{\Sort}^{(2)}_{k,m}$ by $(c,x,y)$ and denote the output by $(a,b)$.
\begin{originalresult}{Claim 11}
There exists a constant $L$ such that any path in $\widetilde{\Sort}^{(2)}_{k,m}$ from an input $\{x_i,y_i\}$ to an output $\{a_j,b_j\}$ has length at most $L\cdot(j-i+1)$, and any path from the input $c$ to an output $\{a_j,b_j\}$ has length at most $L\cdot j$. In particular any path to an output $\{a_j,b_j\}$ has length at most $L\cdot j$.
\end{originalresult}
\begin{proof}
Let $(x',y')$ denote the input to the recursive call to $\widetilde{\Sort}^{(2)}_{k-1,m-1}$ and let $(a',b')$ denote the output of this call. We consider several cases separately:
\begin{itemize}
\item Paths to $\{a_1,b_1\}$ have length bounded by an absolute constant because $a_1$ and $b_1$ are just a (constant-sized) function of the inputs $(c,x_1,y_1)$. If $L$ is sufficiently large then this constant is at most $L$.
\item For $j>1$, paths from $\{c,x_1,y_1\}$ to $a_j$ $(=a'_{j-1})$ consist of a path through a constant-sized circuit (the computation of $c'$) concatenated with a path to $a'_{j-1}$ in $\widetilde{\Sort}^{(2)}_{k-1,m-1}$. By induction the latter path has length at most $L\cdot(j-1)$, so if $L$ is sufficiently large then the total path length is at most $L\cdot j$.
\item For $i>1$ and $j>1$, paths from $\{x_i,y_i\}$ to $\{a_j,b_j\}$ are just paths from $\{x'_{i-1},y'_{i-1}\}$ to $\{a'_{j-1},b'_{j-1}\}$ and thus by induction they have length at most $L\cdot(j-i+1)$.
\end{itemize}
\end{proof}
Returning to $\Sort^{(2)}_{k,m}=\widetilde{\Sort}^{(2)}_{k,m}(?,\cdot,\cdot)$ (without the tilde), we can see that any input-to-output path in $\Sort^{(2)}_{k,m}$ directly corresponds to an input-to-output path in $\widetilde{\Sort}^{(2)}_{k,m}$ whose length is appropriately bounded by Claim 11. That concludes the proof of Lemma 10.
\end{proof}
